\subsection{Transform assumptions and Fourier--Dedekind reduction}
Throughout this section, $p\geq2$, $q$ is a unit modulo $p$, the lift
$\widetilde q$ is odd, $u=q^{-1}\pmod p$, $N=J(J+1)$, and $J\geq1$.
We take \eqref{eq:xi-transform} on nonidentity classes and
$\widehat\xi_J(0)=N(2J+1)/12$ as the defining data.  No filling or
unresolved geometric sign choice enters the following proofs.
Set $\alpha=(1+\widetilde q)/2$, $\beta=(1-\widetilde q)/2$ and
\[
 G(x)=x^\alpha\chi_J(x^\alpha)\chi_{J-1}(x^\beta)
     =\sum_e g_e x^e .
\]
The factorization $P_k=x^{-\alpha}(1-x)(1-x^q)$ at $x=\exp(2\pi ik/p)$
gives the exact Fourier inversion
\begin{equation}\label{eq:fd-inversion}
 \xi_{J,a}=\frac{N(2J+1)}{12p}+\sum_e g_e S_{e+a},\qquad
 S_m=\frac1p\sum_{\substack{\lambda^p=1\\\lambda\ne1}}
       \frac{\lambda^m}{(1-\lambda)(1-\lambda^q)}.
\end{equation}
Fourier--Dedekind sums are classical; see \cite{BeckRobins2004}.
Define the ordinary Dedekind sum by
$s(q,p)=\sum_{j=1}^{p-1}(j/p-1/2)((qj\bmod p)/p-1/2)$,
where remainders are least nonnegative integers.  Direct finite summation gives
\begin{align}
 S_m&=\frac1{p^2}\sum_{j=0}^{p-1}j\,((-m-qj)\bmod p)
                    -\frac{(p-1)^2}{4p},\\
 S_0&=-s(q,p)+\frac{p-1}{4p},\qquad
 S_{m+1}-S_m=\frac{(-mu\bmod p)}p-\frac{p-1}{2p}.
\end{align}
For example, the first identity follows from
$(\lambda-1)\sum_{j=0}^{p-1}j\lambda^j=p$ at nonidentity roots and
orthogonality of the $p$th roots.  Summing the difference, also backwards
for negative $m$, yields
\begin{equation}\label{eq:fd-quadratic}
 S_m-S_0\equiv Q(m):=-\frac{u m(m-1)+(p-1)m}{2p}\pmod{\Z}.
\end{equation}

\subsection{The phase law and rank-one generation}
\begin{theorem}[All-rank phase formula]\label{thm:phase-formula}
For the stated transform, every rank and flat label satisfy
\begin{align}
 \xi_{J,a}-\xi_{J,0}&\equiv-\frac{Nu a^2}{2p}\pmod{\Z},
                                      \label{eq:phase-law}\\
 \xi_{J,a}&\equiv-Ns(q,p)-\frac{Nu a^2}{2p}
 -\frac{N(N-2)(\widetilde q+u)}{12p}-\frac N4\pmod{\Z}.
                                      \label{eq:phase-closed}
\end{align}
Both spin lifts give the same phases.
\end{theorem}
\begin{proof}
The first three moments of $G$ are
\begin{align*}
 G_0&:=\sum_e g_e=N,\qquad G_1:=\sum_e eg_e=\alpha N,\\
 G_2&:=\sum_e e^2g_e=\alpha^2N
 +\frac{\alpha^2J^2(J+1)(J+2)+\beta^2J(J+1)^2(J-1)}3.
\end{align*}
The character's first logarithmic derivative at one vanishes; its second
is $\gamma^2r(r+1)(r+2)/3$ for $\chi_r(x^\gamma)$.
Equations~\eqref{eq:fd-inversion}--\eqref{eq:fd-quadratic} give the relative
phase as $-Nu a^2/(2p)$ plus
$-aN(u\widetilde q-1+p)/(2p)$, which is integral.
For the absolute phase, the same equations give
\[
 \xi_{J,0}\equiv\frac{N(2J+1)}{12p}+NS_0
 -\frac{u(G_2-G_1)+(p-1)G_1}{2p}\pmod{\Z}.
\]
Put $t=\widetilde q$ and $ut=1+pk$. Substitution reduces the difference
between this expression and the right side of \eqref{eq:phase-closed}
at $a=0$ to integers: the nontrivial term is
$Nk[t(N+1)+2J+1]/12$.
Indeed $12\mid N[t(N+1)+2J+1]$: modulo three either $3\mid N$ or
$J\equiv1$ and the bracket is zero; modulo four either $4\mid N$ or
$N\equiv2$ and the bracket is even.  Also $N(t-1)/4$ is integral.
Finally $24\mid N(N-2)=(J-1)J(J+1)(J+2)$, so changes of lift or of the
representative of $u$ change the formula by integers.
\end{proof}

The classical congruence
$12p\,s(q,p)\equiv q+u\pmod p$, including the integrality of $12p\,s$,
is an arithmetic property of Dedekind sums \cite[Ch.~3]{RademacherGrosswald1972}.
It simplifies \eqref{eq:phase-closed} to
\begin{equation}\label{eq:phase-simple}
 \xi_{J,a}\equiv-N(N-1)s(q,p)-\frac{Nu a^2}{2p}-\frac N4\pmod{\Z}.
\end{equation}
\begin{corollary}[Rank one generates all phases]\label{cor:rank-one}
With $c_J=J(J+1)/2$ and $\tau_{J,a}=\exp(2\pi i\xi_{J,a})$,
\begin{align}
 \xi_{J,a}&\equiv c_J\xi_{1,a}+2c_J(c_J-1)\xi_{1,0}\pmod{\Z},\\
 \tau_{J,a}&=\tau_{1,a}^{c_J}\tau_{1,0}^{\,2c_J(c_J-1)}.
\end{align}
Thus the complete exponentiated rank ladder carries no more information
than the labelled rank-one phases, under the same coherent relabeling.
\end{corollary}
\begin{proof}
Insert $\xi_{1,a}\equiv-2s-u a^2/p-1/2$ in \eqref{eq:phase-simple}.
The remaining constant difference is $N(N-2)/4\in\Z$.
\end{proof}
This does not identify the full real-valued corrections, nor eliminate the
information in their integer lifts.

\subsection{A complete phase-collapse criterion}
\begin{theorem}[Coherent-label collapse]\label{thm:phase-collapse}
For the same modulus $p$, units $q,q',n$, and any allowed lifts,
$\xi_{J,a}(p,q)\equiv\xi_{J,na}(p,q')\pmod\Z$ for every $J,a$ if and only if
\[
 q'\equiv n^2q\pmod p,\qquad
 2s(q,p)\equiv2s(q',p)\pmod{\Z}.
\]
\end{theorem}
\begin{proof}
At rank one, label zero imposes the Dedekind condition. Subtracting that
case from label one imposes $u\equiv n^2u'\pmod p$, equivalent to the
first condition. Conversely these conditions give equality at every label
of rank one; Corollary~\ref{cor:rank-one} gives every rank.
\end{proof}
The first condition is the usual oriented homotopy relation, weaker than
oriented homeomorphism.  The second is additional arithmetic information.
Equality questions modulo integers for Dedekind sums have an established
literature \cite{Tsukerman2014}; the theorem is an application to the
specified transform, not a new theory of Dedekind sums.

\begin{corollary}[The $p=16$ example]
The pair $L(16,1)$, $L(16,9)$ is not orientation-preservingly homeomorphic,
but has identical phase ladders under $a\mapsto5a$ (also $3a,11a,13a$).
\end{corollary}
\begin{proof}
The oriented homeomorphism criterion fails since $9\ne1^{\pm1}$ modulo 16.
But $5^2\equiv9$ and $12\cdot16\,s(1,16)=210$,
$12\cdot16\,s(9,16)=18$; their difference is $192$, divisible by $6p$.
Apply Theorem~\ref{thm:phase-collapse}.
\end{proof}
The separate quadratic-in-rank calculation
and its 1,536 integer coefficients remain in the supplement as an independent
finite certificate of this example, not the proof of the general theorem.

\subsection{Infinite blind adjacent-rank families}
The following statements concern the exact $\xi$ values, not just phases.
\begin{theorem}[Blind families, including composite moduli]\label{thm:blind}
Let $m\geq5$ be odd and $\gcd(q,m)=\gcd(q^2-1,m)=1$.
For $p=m$ and $r\in\{m-1,m\}$,
\[
 \xi_{r,a}=\frac{r(r+1)(2r+1)}{12p}\qquad\hbox{for all }a.
\]
For $p=2m$, take $q$ odd with the same conditions modulo $m$. Then
\[
 \xi_{r,a}=\frac1p\left\{\frac{r(r+1)(2r+1)}{12}
       +(-1)^{r-1+a}\frac{r(r+1)}4\right\},\quad r=m-1,m.
\]
In either family the two rank vectors are independent of the allowed $q$.
\end{theorem}
\begin{proof}
Writing $z=\zeta^{e_L}$ and $w=\zeta^{e_R}$, centrality means
$p\mid k(1+\widetilde q)$ or $p\mid k(1-\widetilde q)$.
The gcd hypotheses leave only $k=0$ for $p=m$, and $k=0,m$ for $p=2m$.
Since $p\mid2m$, $\chi_{m-1}(z)=0$ whenever $z\ne\pm1$,
and likewise for $w$. Thus the reduced transforms at both ranks vanish
on all noncentral classes. For $p=m$ only the identity term remains.
For $p=2m$, at $k=m$ one has $w=-z$; substitution in the character
polynomials gives $\widehat\xi_r(m)=(-1)^{r-1}r(r+1)/4$.
Fourier inversion multiplies this by $(-1)^a/p$, as asserted.
\end{proof}
For odd $m\geq5$, an admissible $q$ exists exactly when $3\nmid m$.
If $3\mid m$, every unit has $q^2\equiv1\pmod3$, contradicting the
hypothesis. Conversely, $q=2$ works when $3\nmid m$; for the $p=2m$
family, its odd lift $q=m+2$ works. Primality is not necessary: the
hypotheses hold at $m=25,35,49$ with $q=2$, and with $q=3$ at $m=25,35$.
The prime subfamily includes infinitely many nonhomeomorphic collisions:
for primes $p\geq11$, the $p-3$ allowed residues cannot all lie in one
unoriented homeomorphism orbit, whose size is at most four.
The partial-blindness cases $p=3(J+1)$ lie outside these fully blind
families; their classification and a general positive exact-value
reconstruction criterion are not settled by this theorem.

\paragraph{Finite controls, not additional theorem hypotheses.}
The phase formula and rank-one identity each reproduce all 23,688 supplied
table values. The quotient-ring route also checks 3,644 selected prime and
twice-prime entries and 1,690 composite-family entries. These are
reproduced controls, not prospective predictions. Representative-pair
censuses distinguish orientation-preserving from unoriented equivalence;
they must not be read as counts after quotienting every homeomorphism orbit.
